% Options for packages loaded elsewhere
\PassOptionsToPackage{unicode}{hyperref}
\PassOptionsToPackage{hyphens}{url}
\documentclass[
  ignorenonframetext,
]{beamer}
\newif\ifbibliography
\usepackage{pgfpages}
\setbeamertemplate{caption}[numbered]
\setbeamertemplate{caption label separator}{: }
\setbeamercolor{caption name}{fg=normal text.fg}
\beamertemplatenavigationsymbolsempty
% remove section numbering
\setbeamertemplate{part page}{
  \centering
  \begin{beamercolorbox}[sep=16pt,center]{part title}
    \usebeamerfont{part title}\insertpart\par
  \end{beamercolorbox}
}
\setbeamertemplate{section page}{
  \centering
  \begin{beamercolorbox}[sep=12pt,center]{section title}
    \usebeamerfont{section title}\insertsection\par
  \end{beamercolorbox}
}
\setbeamertemplate{subsection page}{
  \centering
  \begin{beamercolorbox}[sep=8pt,center]{subsection title}
    \usebeamerfont{subsection title}\insertsubsection\par
  \end{beamercolorbox}
}
% Prevent slide breaks in the middle of a paragraph
\widowpenalties 1 10000
\raggedbottom
\AtBeginPart{
  \frame{\partpage}
}
\AtBeginSection{
  \ifbibliography
  \else
    \frame{\sectionpage}
  \fi
}
\AtBeginSubsection{
  \frame{\subsectionpage}
}
\usepackage{iftex}
\ifPDFTeX
  \usepackage[T1]{fontenc}
  \usepackage[utf8]{inputenc}
  \usepackage{textcomp} % provide euro and other symbols
\else % if luatex or xetex
  \usepackage{unicode-math} % this also loads fontspec
  \defaultfontfeatures{Scale=MatchLowercase}
  \defaultfontfeatures[\rmfamily]{Ligatures=TeX,Scale=1}
\fi
\usepackage{lmodern}
\ifPDFTeX\else
  % xetex/luatex font selection
\fi
% Use upquote if available, for straight quotes in verbatim environments
\IfFileExists{upquote.sty}{\usepackage{upquote}}{}
\IfFileExists{microtype.sty}{% use microtype if available
  \usepackage[]{microtype}
  \UseMicrotypeSet[protrusion]{basicmath} % disable protrusion for tt fonts
}{}
\makeatletter
\@ifundefined{KOMAClassName}{% if non-KOMA class
  \IfFileExists{parskip.sty}{%
    \usepackage{parskip}
  }{% else
    \setlength{\parindent}{0pt}
    \setlength{\parskip}{6pt plus 2pt minus 1pt}}
}{% if KOMA class
  \KOMAoptions{parskip=half}}
\makeatother
\setlength{\emergencystretch}{3em} % prevent overfull lines
\providecommand{\tightlist}{%
  \setlength{\itemsep}{0pt}\setlength{\parskip}{0pt}}
% Slide layout defaults for warning-clean scientific decks.
\usepackage{etoolbox}
\IfFileExists{xurl.sty}{\usepackage{xurl}}{}
\IfFileExists{seqsplit.sty}{\usepackage{seqsplit}}{\newcommand{\seqsplit}[1]{#1}}
\protected\def\breakseq#1{\seqsplit{#1}}
\protected\def\breaktt#1{\begingroup\ttfamily\seqsplit{#1}\endgroup}
\setlength{\emergencystretch}{6em}
\tolerance=5000
\hbadness=10000
\hfuzz=1pt
\setlength{\tabcolsep}{2pt}
\AtBeginEnvironment{longtable}{\tiny\renewcommand{\arraystretch}{0.86}\setlength{\tabcolsep}{1pt}}
\AtBeginEnvironment{tabular}{\tiny\renewcommand{\arraystretch}{0.86}\setlength{\tabcolsep}{1pt}}

% Natbib and cross-reference fallbacks — slides don't load natbib
% or cleveref, but combined-PDF manuscript prose may emit these
% commands. The fallback renders citations as a bracketed key list
% and cross-references as detokenized labels so slides don't fail on
% undefined control sequences or raw underscores. \providecommand is
% a no-op if packages load later.
\providecommand{\citep}[1]{[#1]}
\providecommand{\citet}[1]{#1}
\providecommand{\citealp}[1]{#1}
\providecommand{\citeauthor}[1]{#1}
\providecommand{\citeyear}[1]{#1}
\providecommand{\cref}[1]{\texttt{\detokenize{#1}}}
\providecommand{\Cref}[1]{\texttt{\detokenize{#1}}}

\newtheorem{proposition}{Proposition}
\newtheorem{hypothesis}{Hypothesis}
\usepackage{bookmark}
\IfFileExists{xurl.sty}{\usepackage{xurl}}{} % add URL line breaks if available
\urlstyle{same}
% fallback for those not using the hyperref driver hyperxmp:
\makeatletter
\@ifundefined{xmpquote}{\newcommand{\xmpquote}[1]{#1}}{}
\makeatother
\hypersetup{
  hidelinks,
  pdfcreator={LaTeX via pandoc}}

\author{\texorpdfstring{}{}}
\date{}

\begin{document}

\section{Introduction}\label{sec:introduction}

\begin{frame}[allowframebreaks]{Introduction}
A \emph{data descriptor} (or data paper) publishes a dataset as a
citable research object. Unlike an analysis paper, its contribution is
not a statistical finding but a \textbf{contract}: a precise,
machine-readable account of what the dataset contains, how it was
produced, how it is licensed, and what quality guarantees it carries.
The FAIR guiding principles --- Findable, Accessible, Interoperable,
Reusable \breakseq{[@wilkinson2016fair]} --- frame why this matters: a dataset
is only reusable to the extent that its structure and provenance are
legible to both humans and machines.
\end{frame}

\begin{frame}[allowframebreaks]{What this template provides}
\protect\phantomsection\label{what-this-template-provides}
This exemplar packages the recurring moving parts of a data descriptor
so that a fork can start from a working, tested baseline:

\begin{itemize}
\tightlist
\item
  A \textbf{schema / data dictionary}: named fields with declared types,
  nullability, units, and value constraints (patterns, enumerations,
  numeric bounds).
\item
  A \textbf{file inventory}: each shipped file with its media type, a
  sha256 checksum, and a row count.
\item
  A \textbf{provenance chain}: the ordered steps and agents that
  produced the release.
\item
  A \textbf{license boundary}: an explicit reuse license, and the
  discipline of publishing \emph{metadata and checksums} rather than
  bundling restricted bytes.
\item
  A \textbf{validation gate}: a tested library that rejects malformed,
  unsafe, or incomplete descriptors and verifies declared checksums
  against real bytes.
\end{itemize}
\end{frame}

\begin{frame}[allowframebreaks]{Scope and honesty}
\protect\phantomsection\label{scope-and-honesty}
The shipped dataset is \textbf{synthetic and deliberately small} --- it
exists to make the workflow runnable and testable, not to support any
empirical claim. Its values are generated deterministically.
Accordingly, this manuscript restricts its claims to dataset
\emph{structure, provenance, quality, and release readiness}. When a
real dataset is forked in, the same validation and figure code applies
unchanged; only the descriptor and fixture files are replaced. The
sections that follow describe the demonstration dataset
(\breakseq{[@sec:data]}), its schema (\breakseq{[@sec:schema]}), its provenance
(\breakseq{[@sec:provenance]}), and the quality gate that binds the descriptor
to the bytes it names (\breakseq{[@sec:quality]}), and close with usage and
forking notes (\breakseq{[@sec:usage]}).
\end{frame}

\end{document}
